\phantomsection\bheading{section}{1.\quad The question, restated precisely}{section}{\protect\numberline{1.}The question, restated precisely}

\phantomsection\bheading{subsection}{1.1\quad What the owner asked}{subsection}{\protect\numberline{1.1}What the owner asked}

“Editing a \code{List} while anything else still holds the old version is a compile error.” Unpacked:

\begin{itemize}
\item \emph{Editing} is a call of one of \code{core/List}'s update primitives: today \code{set}, \code{update}, \code{push}, \code{pop}, \code{swap}, \code{insertAt}, \code{removeAt}, and the two the syntax lowers to, \code{cons} (\code{[ x, …xs ]}) and \code{append} (\code{[ …xs, …ys ]}, \code{++} on lists). §10 asks whether \code{map}, \code{filter}, \code{indexedMap}, \code{reverse}, \code{sort} and \code{take} should join them.
\item \emph{Anything else} is any holder of the cell, of any kind listed in §2.2 — the operand's own binding included, if it is read again after the edit.
\item \emph{Still holds} means the holder is live: some execution continuing from the edit reads through it.
\item \emph{The old version} is the cell: the very array, not an equal list.
\item \emph{A compile error}, not a copy and not a warning. The owner wants no silent fallback: an accepted program writes in place, always, in development and release alike.
\end{itemize}

So the rule is \textbf{uniqueness at the update site, decided by liveness}, with the one escape hatch \code{List.copy} making intentional sharing explicit.

\phantomsection\bheading{subsection}{1.2\quad What “observe” means, and what it does not}{subsection}{\protect\numberline{1.2}What “observe” means, and what it does not}

An in-place write is observable when some computation reads the cell's elements after the write through a reference it obtained before. Three readers of a cell exist in beni:

\begin{enumerate}
\item \textbf{Program code} — any expression that reads elements: \code{List.get}, a pattern, a fold, \code{==} (\code{List.eq} walks both lists), \code{Debug.toString}, a markup hole showing the list.
\item \textbf{The runtime} — today's \code{browser} keeps the rendered list to skip the walk when the next one is identical (\code{forPosition}'s \code{s.b}, \code{backend.md} §15.5); \code{browser-direct} keeps instance items and leaf slots (§5.3, §6.1).
\item \textbf{JavaScript} — a sibling or a \code{Js.from} that kept the array.
\end{enumerate}

Each is a content holder. An \textbf{identity holder} is different: it keeps the reference only to ask \code{===} later. In-place writes do not change identity, so an identity holder never sees a different \emph{object}; what it may do is conclude “unchanged” from “same object”, which is wrong after an in-place write. That hazard is the renderer's, not the program's: beni has no reference equality a program can call (\code{language.md} §11.12), and \code{browser-direct} §7.4's first rule says a group “never decides 'unchanged' by the identity of an object on an in-place-written path”. The soundness claim of §4 covers content holders; identity holders are an assumption on the platforms (A4), discharged by that rule.

\phantomsection\bheading{subsection}{1.3\quad Two semantics that must coincide}{subsection}{\protect\numberline{1.3}Two semantics that must coincide}

The precise form of the guarantee is Cogent's, as O'Connor, Linares Arévalo and Rizkallah put it: uniqueness works “by statically ruling out every program where the difference between the immutable and the mutable interpretations can be observed” (2026, vimpl.tex:45). Aspinall, Hofmann and Konečný prove exactly such a theorem for their first-order system: an in-place evaluation and a “safe” functional evaluation agree on every well-typed term (2008 Thm 5.8, p. 30, “in-place update evaluation is correct … and complete”). §4 states beni's version: for an accepted program, the \emph{value semantics} (every update copies) and the \emph{update semantics} (every accepted update writes the cell) produce the same observations — the same printed output, the same DOM transcript, the same sequence of calls into JavaScript with the same marshalled arguments.

\phantomsection\bheading{subsection}{1.4\quad What is not asked}{subsection}{\protect\numberline{1.4}What is not asked}

\begin{itemize}
\item Records, tuples and constructors are not written in place by this checker. \code{browser-direct} §7.4 (S7) plans record paths; the machinery here extends to them (§7.7) but the thesis is about lists, as the owner asked.
\item The checker does not decide \emph{whether} an update is worth doing in place; every accepted update is in place. There is no cost model and no “copy if small”.
\item It does not replace the write-set analysis; it reads one fact from it (§5.9).
\end{itemize}
